\par\noindent\textbf{(1)} 取场 \(\propto e^{i(\beta x-\omega t)}e^{-\kappa_dz}\) 在介质侧、\(e^{+\kappa_mz}\) 在金属侧，\(\kappa_j^2=\beta^2-\varepsilon_jk_0^2\)。TM模式具有 \(E_x,E_z,H_y\) 分量，要求 \(\operatorname{Re}\kappa_j>0\) 保证表面局域化。
\par\noindent\textbf{(2)} 无自由面电流时切向 \(E_x,H_y\) 连续。使用TM波阻抗关系可化为 \(\kappa_d/\varepsilon_d+\kappa_m/\varepsilon_m=0\)。联立各自 \(\kappa_j^2\)，得到 \(\beta=k_0\sqrt{\varepsilon_d\varepsilon_m/(\varepsilon_d+\varepsilon_m)}\)，在所给符号范围内为实大于介质光线。
\par\noindent\textbf{(3)} 每侧场幅 \(1/e\) 穿透长度为 \(\ell_j=1/\kappa_j\)。当 \(\varepsilon_m\to-\varepsilon_d^-\) 时分母 \(\varepsilon_d+\varepsilon_m\to0^-\)，理想局域波矢 \(\beta\to\infty\)、局域长度趋零，意味着准静态表面等离激元极限。
\par\noindent\textbf{(4)} 真实金属 \(\varepsilon_m(\omega)\) 带虚部，\(\beta\) 变复数并使沿界面传播有限；当 \(\beta\) 接近电子微观非局域尺度时局域Drude近似失效。不能把理想无限大波矢解释为无需能量、无限分辨率的可实现模式。
